\documentclass[12pt]{article}

% The preamble file contains all the macros and page setup. I recommend making a folder where you compile your exercises and always call on the preamble.
%%%%%%%%%%%%%%%%%%%%%%%%%%%%%%%%%%%%%%%%%%%%%%%%%%%%%%%%%%%%%%%%%%%%%%
%% all the formatting stuff and packages



\usepackage{geometry}
\geometry{
  top=4cm,
  headheight=70pt,
  headsep=25pt,
  inner=1.25in,
  outer=1.25in
}

\usepackage[utf8]{inputenc}

\usepackage{xcolor,graphicx,hyperref,amsthm,amsmath,amssymb}
\usepackage{algpseudocode}


\usepackage{fancyhdr}
\pagestyle{fancy}

%%% Begin Header Definitions %%%
\fancypagestyle{firststyle}
{
   	\lhead{}
	\rhead{\raisebox{0.2cm}{\includegraphics[width=3.5cm]{resource/hpi_logo.pdf}}}
	\renewcommand{\headrulewidth}{0pt}
}


\newcommand{\makeheader}[1]{
\noindent
\textbf{Algorithmix \hfill Winter term 2019/20 \\ Andreas G\"obel}

\begin{center}
\large{\textbf{\makeHeadline}}\\
\end{center}

\begin{center}
\textit{#1} 
\end{center}


}

\newcommand{\lecturer}[1]{\newcommand{\makeLecturer}{#1}}
\newcommand{\lectureName}[1]{\newcommand{\makeLectureName}{#1}}
\newcommand{\semester}[1]{\newcommand{\makeSemester}{#1}}
\newcommand{\headline}[1]{\newcommand{\makeHeadline}{#1}}
\newcommand{\urlLine}[1]{\newcommand{\makeUrlLine}{#1}}
%%% End Header Definitions %%%

\newcommand{\pr}{\ensuremath{\mathbb{P}}}
\newcommand{\E}{\ensuremath{\mathbb{E}}}


\newcommand{\makeCompact}{%
\setlength{\itemsep}{1pt}%
\setlength{\parsep}{0pt}%
\setlength{\parskip}{0pt}}


\newcommand{\Var}{\mathrm{Var}}


\newcommand{\student}[2]{
\thispagestyle{firststyle}

%%% Begin Header Specification %%%

\headline{Solutions to homework #1}

\makeheader{#2}

\medskip

}

\newcommand{\sol}[1]{ \hspace{.5cm}\\ \noindent \textbf{Exercise #1.}}

% Here you can define your own macros
\newcommand{\calF}{\ensuremath{\mathcal{F}}}
\newcommand{\polyreduce}{$\leq_{m}^{p}$}

\begin{document}
 
 %To create a header for the document use the following macro. \student{Homework number}{Name}

 \student{5}{Stefan Reschke}
 

 \sol{1} The Steiner tree problem (STP) can indeed be approximated to this extend. We write MStT for the minimum Steiner tree to not confuse it with the minimum spanning tree (MST).
 
 Since the MST spans all vertices $V$, it also spans all vertices $V_0$, since $V_0\subseteq V$.
 
 We consider a DFS traversal of $G$, including every return to a vertex at the end of a subtree in our recursive calls. This traversal is a cycle over all $v\in V$. Let that traversal be of length $m$, we have \begin{equation}
 	\sum_{i=0}^{m} d(v_i, v_i+1) = 2\cdot \texttt{cost}(\texttt{MStT}),
 \end{equation} since all edges are used twice. We remove all vertices in $V\setminus V_0$ and only keep the first occurences of all in $V_0$. We say this new traversal is of length $k$. Then we have \begin{equation}
 	\sum_{i=0}^{k} d(v_i, v_i+1) \leq \sum_{i=0}^{m} d(v_i, v_i+1)
 \end{equation}, since $k<m$a and thus \begin{equation}
 	\sum_{i=0}^{k} d(v_i, v_i+1) \leq 2\cdot \texttt{cost}(\texttt{MStT}).
 \end{equation} This $k$-sized traversal being a path, is also a tree. Using the graph induced, we arrive at the assumption.

 \sol{2} (a) Both operations on RB-Trees can involve a rotation. However, they don't have to. Possible outcomes are: \begin{enumerate}
 \item \texttt{insert} requires rotation, \texttt{deletion} requires rotation
 \item \texttt{insert} requires no rotation, \texttt{deletion} requires rotation
 \item \texttt{insert} requires rotation, \texttt{deletion} requires no rotation
 \item \texttt{insert} requires no rotation, \texttt{deletion} requires no rotation
 \end{enumerate}

\vspace{.5cm}
\noindent (b)
 
  
\end{document}
